% ============================================================================
\chapter{Static Arbitrage and Parametric Surfaces}\label{ch:theory}
% ============================================================================

A smoother is judged on two things it cannot optimise simultaneously: how closely it follows the
quotes, and whether the prices it manufactures between them are coherent. This chapter builds
the second half of that sentence. It answers three questions in order, and each answer is used
somewhere later in the thesis.

\emph{What exactly must a surface avoid?} Section~\ref{sec:setting} defines static arbitrage,
shows why the abstract definition reduces to two concrete conditions, one in the strike
direction and one in the maturity direction, and explains why those two exhaust the problem
rather than merely being convenient. \emph{How do we test for it?}
Sections~\ref{sec:blacksec} to~\ref{sec:calendarsec} translate both conditions into the
total-variance coordinates $(k,w)$, where the butterfly condition collapses to the sign of a
single scalar function $g$ and the calendar condition to the sign of a single derivative
$\partial_\tau w$. These are the two numbers audited on every surface in this thesis.
\emph{Why should anyone care?} Section~\ref{sec:dupiresec} derives Dupire's identity
$\sigma^2_{\mathrm{loc}}=\partial_\tau w/g$, which turns out to say that the two audited
quantities are the numerator and the denominator of the object desks actually price with. That
identity is what allows Chapter~\ref{ch:downstream} to convert a violation rate into an economic
cost, and it is the reason the audit is worth performing at all.

Sections~\ref{sec:svisec} and~\ref{sec:ssvisec} then introduce the parametric families, SVI and
SSVI, for which the conditions of this chapter are not merely testable but \emph{provable}. They
serve three roles later: as the benchmark every learned model must beat, as the certified prior
inside the deep smoother, and as the source of the closed forms against which our automatic
differentiation is validated. Chapter~\ref{ch:learning} then turns to model classes for which no
such proof exists.

\paragraph{Attribution.} Because this chapter mixes established theory with results developed
for this thesis, the division is stated once here rather than left implicit. The
static-arbitrage framework, the Black partials, the Breeden--Litzenberger identity, Dupire's
equation, the SVI parameterisations and the Gatheral--Jacquier certificates are due to the cited
authors, and every statement carries its source; our contribution to these sections is
verification rather than derivation, since each result is re-derived independently and checked
numerically against the values published by its authors. Two results are exceptions and are
flagged in place: the correction to the published inverse of the natural-parameter map
(Remark~\ref{rem:gjtypo}), and the endpoint-certification result (Theorem~\ref{lem:range}),
which reduces the Gatheral--Jacquier conditions on a compact working range to two endpoint
checks and is what allows a per-day prior to be \emph{certified} rather than merely audited.
Section~\ref{sec:softtheory} of the next chapter is original throughout. Proofs of published
results are reproduced when they are short and used later, and cited when they are long; every
compression is signalled explicitly.

To make the division visible rather than merely stated, the thesis adopts a naming convention
throughout: results established by other authors are labelled \emph{Propositions} or
\emph{Lemmas} and carry an explicit citation in their statement, whereas \emph{Theorems} are
reserved for results proved here. There are exactly two of the latter,
Theorem~\ref{lem:range} in this chapter and Theorem~\ref{prop:nodegap} in
Chapter~\ref{ch:learning}.

% ============================================================================
\section{Market setting and static arbitrage}\label{sec:setting}
% ============================================================================

Fix a trading date $t$ and let $S=(S_u)_{u\ge t}$ denote the underlying index. For a maturity
$T>t$ write $\tau=T-t$, let $F=F(t,T)$ be the forward price for delivery at $T$ and
$D=D(t,T)$ the discount factor. European calls and puts of strike $K$ and maturity $T$ trade at
prices $C(K,T)$ and $P(K,T)$.

We work throughout under the standard \emph{static} and model-free framework: no dynamic model is
assumed. The only structure imposed on the quoted (or fitted) price surface is the absence of
\emph{static arbitrage}, i.e.\ the impossibility of forming at time $t$, a buy-and-hold
portfolio of finitely many calls, puts, forwards and cash whose cost is negative while its
payoff at the relevant maturities is almost surely nonnegative.

\begin{definition}[Static arbitrage]\label{def:staticarb}
A call price surface $(K,T)\mapsto C(K,T)$ admits \emph{static arbitrage} if there exist
$n\in\mathbb N$, strikes $K_1,\dots,K_n$, maturities $T_1,\dots,T_n$, positions
$\alpha_1,\dots,\alpha_n\in\R$ in the corresponding calls, positions in forwards and cash, such
that the portfolio's initial cost is strictly negative while its terminal payoff is nonnegative
in every state. The surface is \emph{free of static arbitrage} if no such portfolio exists.
\end{definition}

Definition~\ref{def:staticarb} is a statement about \emph{portfolios}, whereas everything this
thesis fits and audits is a \emph{surface}. The bridge between the two deserves to be made
explicit at the outset, because it is what turns an abstract condition into something a smoother
can actually violate. A fitted surface is not a drawing. Through the Black formula it defines a
complete system of option prices,
\[
w(k,\tau)\;\longrightarrow\;\sigma_{\mathrm{IV}}(k,\tau)\;\longrightarrow\;C(K,T),
\]
including at strikes and maturities where nothing was ever quoted. Those manufactured prices
must satisfy the same inequalities as the quoted ones, and an interpolation that honours every
observation point by point can still generate a system that does not.

Two failures are possible, and they correspond to the two directions of the surface. In the
\emph{strike} direction, fix a maturity and take three strikes $K_1<K_2<K_3$ with
$K_2=\lambda K_1+(1-\lambda)K_3$. Convexity of the call price in $K$ requires
\[
C(K_2)\;\le\;\lambda\,C(K_1)+(1-\lambda)\,C(K_3),
\]
and if the fitted surface breaks this inequality then the butterfly portfolio that is long the
wings and short the body has a nonnegative payoff and a negative cost: free money. In the
\emph{maturity} direction, a longer-dated option must be worth at least as much as a
shorter-dated one at the same moneyness, which in total-variance coordinates will read
$\partial_\tau w\ge0$; otherwise the calendar spread that is long the far leg and short the near
leg has the same defect. The audit performed throughout this thesis is nothing more than
checking these two families of inequalities on the surfaces our models produce, and the whole of
Chapter~\ref{ch:downstream} is the question of what it costs when they fail.

It is natural to ask why these two directions and not some longer list. The answer, and it is
what makes a two-constraint audit exhaustive rather than merely convenient, is that under the
standard positivity, regularity and boundary conditions, which are largely structural (total
variance positive, correct behaviour as $\tau\downarrow0$, the price bounds of
Proposition~\ref{prop:shape}, and the wing behaviour of Remark~\ref{rem:lee}), the two
\emph{nontrivial} directions of static arbitrage are exactly these:
\[
\text{no static arbitrage}
\iff
\begin{cases}
\text{no butterfly arbitrage on every maturity slice,}\\[2pt]
\text{no calendar-spread arbitrage across maturities.}
\end{cases}
\]
This decomposition is Lemma~2.1 of \cite{gatheraljacquier2014}, resting on
\cite{coxhobson2005,kellerer1972}. Its sufficiency has a probabilistic reading worth stating
once, because it explains why the two conditions are the right ones rather than two conditions
among many. A butterfly-free slice is precisely a slice whose implied density is a genuine
probability distribution, so each maturity contributes an admissible marginal law. The calendar
condition is precisely the statement that these marginals are correctly ordered, in convex
order. And a family of marginals increasing in convex order is, by Kellerer's theorem
\cite{kellerer1972}, the marginal system of some martingale, which is in turn the absence of
static arbitrage. Butterfly therefore governs coherence \emph{within} a maturity, calendar
governs coherence \emph{across} maturities, and together they close the problem.

\newpage

The classical characterisation, which we take as the organising principle of the thesis, is that
freedom from static arbitrage is equivalent to the existence of a nonnegative
\emph{martingale measure} consistent with all quotes: there is a filtered probability space
carrying a nonnegative martingale $(M_u)$ with $M_t=1$ such that
\begin{equation}\label{eq:pricingrule}
C(K,T)=D(t,T)\,F(t,T)\;\E\Bigl[\bigl(M_T-\tfrac{K}{F(t,T)}\bigr)^+\Bigr],
\end{equation}
i.e.\ $M_T=S_T/F(t,T)$ is the maturity-$T$ forward-normalised underlying. Sufficiency of
\eqref{eq:pricingrule} is immediate (linearity and positivity of expectations); necessity is the
content of Kellerer's theorem \cite{kellerer1972} and its refinements for call surfaces
\cite{coxhobson2005,roper2010}: a system of call prices that is (per maturity) decreasing and
convex in $K$ with the correct bounds, and (across maturities) increasing in the sense of
Section~\ref{sec:calendarsec}, is the marginal-expectation system of some martingale. We use
this equivalence freely and now record the elementary per-maturity properties.

\begin{proposition}[Shape of an arbitrage-free call slice]\label{prop:shape}
Fix $T$ and suppose $C(\cdot,T)$ is free of static arbitrage. Write $c(K)=C(K,T)/D(t,T)$ for
the \emph{undiscounted} call, which strips out the interest-rate factor and equals the pure
expected payoff. Then:
\begin{enumerate}[label=(\alph*)]
\item $K\mapsto c(K)$ is nonincreasing and convex;
\item $(F-K)^+\le c(K)\le F$ for all $K\ge0$;
\item if $c$ is twice differentiable, $\partial_{KK}c(K)\ge0$, and $\partial_{KK}c(K)$ is the
(undiscounted) risk-neutral density of $S_T$ at $K$ (Breeden--Litzenberger \cite{breeden1978}).
\end{enumerate}
\end{proposition}

\begin{proof}
\begin{enumerate}[label=\textup{(\alph*)}, leftmargin=*]

\item \textbf{Monotonicity.}
For $K_1<K_2$, the call spread consisting of a long call at $K_1$ and a short call at $K_2$
has payoff $(S_T-K_1)^+-(S_T-K_2)^+\ge 0.$
Therefore, its cost must be nonnegative: $c(K_1)-c(K_2)\ge 0.$

\medskip

\textbf{Convexity.} Let $K_1<K_2<K_3$, with $K_2=\lambda K_1+(1-\lambda)K_3.$
The \emph{butterfly} portfolio consisting of a position $+\lambda$ at $K_1$, a position $-1$
at $K_2$, and a position $+(1-\lambda)$ at $K_3$ has payoff
$\lambda(S_T-K_1)^+-(S_T-K_2)^+ +(1-\lambda)(S_T-K_3)^+\ge 0,$
by convexity of the map $K\mapsto (S_T-K)^+$. Hence, its cost satisfies
$\lambda c(K_1)-c(K_2)+(1-\lambda)c(K_3)\ge 0.$ \medskip

\item The upper bound follows from the dominating payoff $S_T$, whose cost is $F$ and which
can be replicated using a forward position together with cash. The lower bound follows from
the dominated payoff $S_T-K$, which can be replicated using a forward position minus cash.
\medskip
\item Under \eqref{eq:pricingrule}, the undiscounted call is $c(K)=\E\bigl[(S_T-K)^+\bigr].$ Differentiating twice under the expectation gives $\partial_{KK}c(K)=E\bigl[\delta_K(S_T)\bigr],$ which is the density of $S_T$ evaluated at $K$. The positivity of this density is the
infinitesimal counterpart of the convexity established in part~(a). Indeed, the butterfly
portfolio with strikes $K-\varepsilon$, $K$, and $K+\varepsilon$ has cost
$\varepsilon^2\partial_{KK}c(K)+o(\varepsilon^2).$
\end{enumerate}
\end{proof}

\newpage

% ============================================================================
\section{The Black model, implied volatility, and the $(k,w)$ coordinates}\label{sec:blacksec}
% ============================================================================

\subsection{The Black formula}

Under the Black (forward) model, $M_T=S_T/F$ is lognormal: $M_T=\exp(\sqrt{v}\,Z-v/2)$ with
$Z\sim\mathcal N(0,1)$ and total variance $v=\sigma^2\tau$. The call price is
\begin{equation}\label{eq:black}
C(K,\tau;\sigma)=D\,\bigl[F\,N(d_+)-K\,N(d_-)\bigr],\qquad
d_\pm=\frac{\log(F/K)}{\sigma\sqrt{\tau}}\pm\frac{\sigma\sqrt{\tau}}{2},
\end{equation}
with $N$ the standard normal cdf.

% \begin{proof}[Derivation of \eqref{eq:black}]
% $C/D=F\,\E[(M_T-e^{k})^+]$ with $k=\log(K/F)$. On the event
% $\{M_T>e^k\}=\{Z>(k+v/2)/\sqrt v\}=\{Z>-d_-\}$ (with $v=\sigma^2\tau$),
% \[
% \E[M_T\mathbf 1_{\{Z>-d_-\}}]
% =\int_{-d_-}^{\infty} e^{\sqrt v z-v/2}\varphi(z)\dd z
% =\int_{-d_-}^{\infty}\varphi(z-\sqrt v)\dd z
% =N(d_-+\sqrt v)=N(d_+),
% \]
% by completing the square, while $\E[e^k\mathbf 1_{\{Z>-d_-\}}]=e^kN(d_-)$. Multiplying by $F$
% and noting $Fe^k=K$ gives \eqref{eq:black}.
% \end{proof}

Since the Black \emph{vega} $\partial C/\partial\sigma=DF\varphi(d_+)\sqrt\tau>0$, \eqref{eq:black} is strictly increasing in $\sigma$ and therefore inverts on the
arbitrage-free price interval: the \emph{implied volatility} $\sigma_{\mathrm{IV}}(K,\tau)$ of a
market quote is the unique root of $C(K,\tau;\sigma)=C^{\mathrm{mkt}}$. The dependence of implied volatility on $(K,\tau)$ describes the volatility smile across strikes and the term structure across maturities, thereby defining the implied volatility surface.

\subsection{Coordinates}

Two changes of variable make the theory tractable. The \emph{log-moneyness} $k=\log(K/F)$
centres the smile on the forward; the \emph{total (implied) variance} is the natural vertical coordinate:
\begin{equation}\label{eq:w}
w(k,\tau)=\sigma_{\mathrm{IV}}^2(k,\tau)\,\tau
\end{equation}
Both static no-arbitrage conditions, the risk-neutral
density and Dupire's formula all take their simplest form in $(k,w)$. Throughout, primes denote
$\partial_k$ and $\dot w$ or $\partial_\tau w$ the maturity derivative. We also write the
normalized, undiscounted forward prices $c(k,w)=C/(DF)$, $p(k,w)=P/(DF)$, so that
\begin{equation}\label{eq:blacknorm}
c(k,w)=N(d_+)-e^{k}N(d_-),\qquad d_\pm(k,w)=-\frac{k}{\sqrt w}\pm\frac{\sqrt w}{2},
\end{equation}
and put--call parity reads $c-p=1-e^{k}$. Note the identity, used repeatedly below,
\begin{equation}\label{eq:phiid}
\varphi(d_+)=e^{k}\,\varphi(d_-),
\end{equation}
which follows from $d_+^2-d_-^2=(d_++d_-)(d_+-d_-)=(-2k/\sqrt w)(\sqrt w)=-2k$.

\subsection{The Black partial derivatives}

All later derivations reduce to the following closed forms; we prove each in full. 

\begin{lemma}[Black partials in $(k,w)$]\label{lem:blackpartials}
Let $c(k,w)$ be as in \eqref{eq:blacknorm} and let subscripts denote partial derivatives at
fixed $(k,w)$. Then
\begin{align}
c_k &= -e^{k}N(d_-), \label{eq:ck}\\
c_w &= \frac{e^{k}\varphi(d_-)}{2\sqrt w}\;=\;\frac{\varphi(d_+)}{2\sqrt w}\;>0, \label{eq:cw}\\
c_{kk}-c_k &= \frac{e^{k}\varphi(d_-)}{\sqrt w}\;=\;2\sqrt w\, c_w, \label{eq:ckk}\\
c_{kw} &= \Bigl(\frac12-\frac{k}{w}\Bigr)\,c_w, \label{eq:ckw}\\
c_{ww} &= \Bigl(-\frac18-\frac{1}{2w}+\frac{k^2}{2w^2}\Bigr)\,c_w. \label{eq:cww}
\end{align}
\end{lemma}


\begin{proof}
Throughout, $\partial_k d_\pm=-1/\sqrt w$ and
$\partial_w d_\pm=\dfrac{k}{2w^{3/2}}\pm\dfrac{1}{4\sqrt w}$.

\begin{enumerate}[label=\textup{(\roman*)}, leftmargin=*]

\item 
$c_k=\varphi(d_+)\partial_kd_+-e^kN(d_-)-e^k\varphi(d_-)\partial_kd_-
=-e^kN(d_-)-\dfrac{1}{\sqrt w}\bigl(\varphi(d_+)-e^k\varphi(d_-)\bigr)
=-e^kN(d_-)$.

\item 
$c_w=\varphi(d_+)\partial_wd_+-e^k\varphi(d_-)\partial_wd_-
=e^k\varphi(d_-)\bigl(\partial_wd_+-\partial_wd_-\bigr)
=\dfrac{e^k\varphi(d_-)}{2\sqrt w}>0$.\\
Since $\partial w/\partial\sigma=2\sigma\tau$, this recovers the classical
Black vega $DF\varphi(d_+)\sqrt\tau$.

\item 
Differentiating \eqref{eq:ck} in $k$ gives
$c_{kk}=-e^kN(d_-)-e^k\varphi(d_-)\partial_kd_-
=c_k+\dfrac{e^k\varphi(d_-)}{\sqrt w}$.

\item 
Differentiating \eqref{eq:ck} in $w$ gives
$c_{kw}=-e^k\varphi(d_-)\partial_wd_-
=-e^k\varphi(d_-)\bigl(\dfrac{k}{2w^{3/2}}-\dfrac{1}{4\sqrt w}\bigr)
=\dfrac{e^k\varphi(d_-)}{2\sqrt w}\bigl(\dfrac12-\dfrac{k}{w}\bigr)
=\bigl(\dfrac12-\dfrac{k}{w}\bigr)c_w$.

\item 
Differentiating \eqref{eq:cw} in $w$ and using $\varphi'(x)=-x\varphi(x)$ gives
\[
c_{ww}
=e^k\left[
\frac{-d_-\varphi(d_-)\partial_wd_-}{2\sqrt w}
-\frac{\varphi(d_-)}{4w^{3/2}}
\right]
=c_w\left[(-d_-)\partial_wd_--\frac{1}{2w}\right].
\]
Since $-d_-=\dfrac{k}{\sqrt w}+\dfrac{\sqrt w}{2}$,then $ (-d_-)\partial_wd_-
=\left(\frac{k}{\sqrt w}+\frac{\sqrt w}{2}\right)
\left(\frac{k}{2w^{3/2}}-\frac{1}{4\sqrt w}\right)
=\frac{k^2}{2w^2}-\frac18$.
\end{enumerate}
\end{proof}


% ============================================================================
\section{Butterfly arbitrage, the Durrleman function, and the implied density}
\label{sec:butterflysec}
% ============================================================================

We now take the two arbitrage directions in turn, starting with the strike direction. At a fixed
maturity the question is whether the fitted smile corresponds to a valid probability
distribution over the terminal underlying. Asking for a valid distribution is asking for a
nonnegative density; by Breeden--Litzenberger that is asking for convexity of the call price in
strike; and in the $(k,w)$ coordinates that condition collapses to a single scalar built from
the slice and its first two derivatives. That scalar is the Durrleman function $g$, and the
entire butterfly audit of this thesis reduces to monitoring its sign. This section builds the
chain and ends with the complete characterisation, which requires one boundary condition in
addition to $g\ge0$.

Following \cite{gatheraljacquier2014,durrleman2010,roper2010}, a surface is free of static
arbitrage if and only if every maturity slice is free of \emph{butterfly} arbitrage and the
family of slices is free of \emph{calendar} arbitrage (this decomposition is
\cite[Lemma~2.1]{gatheraljacquier2014}, resting on \cite{coxhobson2005,kellerer1972}). This
section treats the butterfly direction; Section~\ref{sec:calendarsec} treats the calendar one.

\begin{definition}[Durrleman function and butterfly arbitrage]\label{def:durrleman}
For a $C^2$ total variance slice $w(\cdot)>0$ at fixed maturity, the associated
\emph{Durrleman function} is defined by
\begin{equation}\label{eq:g}
g(k)=\Bigl(1-\frac{k,w'(k)}{2,w(k)}\Bigr)^{2}
-\frac{w'(k)^{2}}{4}\Bigl(\frac{1}{w(k)}+\frac14\Bigr)
+\frac{w''(k)}{2}.
\end{equation}
The slice is \emph{free of butterfly arbitrage} if and only if the associated call prices
$k\mapsto c(k,w(k))$ are convex in strike and vanish as $K\to\infty$.
\end{definition}


The next proposition identifies $g$ with the risk-neutral
density and thereby with the convexity of the call slice. It is the precise statement of the
slogan \emph{a butterfly violation is negative probability mass}.

\begin{proposition}[Implied density in log-moneyness]\label{prop:density}
Let $X_T=\log(S_T/F_T)$ and let $w(\cdot)$ be a $C^2$ slice with implied call prices
$c(k)=c\bigl(k,w(k)\bigr)$ as in \eqref{eq:blacknorm}. The risk-neutral density of $X_T$ is
\begin{equation}\label{eq:density}
q(k)=\frac{g(k)}{\sqrt{2\pi\,w(k)}}\,
\exp\!\Bigl(-\tfrac12 d_-(k)^2\Bigr)
=\frac{g(k)\,\varphi\bigl(d_-(k)\bigr)}{\sqrt{w(k)}},
\qquad d_-(k)=d_-\bigl(k,w(k)\bigr),
\end{equation}
with $\varphi$ the standard normal pdf. In particular $q\ge 0\iff g\ge 0$.
\end{proposition}

\begin{proof}
\noindent\textbf{Step 1: Density in log-moneyness.}
By Proposition~\ref{prop:shape}(c), the undiscounted density of $S_T$ at $K$ is
$\partial_{KK}(C/D)$. Since $K=Fe^k$, then $q(k)=K\,\partial_{KK}(C/D).$
Using $C/D=Fc(k)$ and $\partial_K=K^{-1}\partial_k$,
we obtain: 
$
\partial_{KK}(C/D)
=\frac1K\partial_k\left[\frac{F}{K}\partial_kc\right]
=\frac{F}{K^2}\bigl(\partial_{kk}c-\partial_kc\bigr),
$
hence
\begin{equation}\label{eq:qk}
q(k)=e^{-k}\bigl(\partial_{kk}c-\partial_kc\bigr).
\end{equation}

\noindent\textbf{Step 2: Total derivatives.}
For $c(k)=c(k,w(k))$,
\[
\partial_kc=c_k+c_ww',
\qquad
\partial_{kk}c=c_{kk}+2c_{kw}w'+c_{ww}(w')^2+c_ww''.
\]

\noindent\textbf{Step 3: Substitution.}
Using Lemma~\ref{lem:blackpartials} in \eqref{eq:qk},
\[
q(k)=e^{-k}c_w\left[
2\sqrt w-\frac{2kw'}{w}
+\left(\frac{k^2}{2w^2}-\frac{1}{2w}-\frac18\right)(w')^2+w''
\right].
\]
Since $e^{-k}c_w=\varphi(d_-)/(2\sqrt w)$,
\[
q(k)=\frac{\varphi(d_-)}{\sqrt w}
\left[
1-\frac{kw'}{w}
+\left(\frac{k^2}{4w^2}-\frac{1}{4w}-\frac{1}{16}\right)(w')^2
+\frac{w''}{2}
\right].
\]
The bracket is precisely $g(k)$ from \eqref{eq:g}, so
\[
q(k)=g(k)\frac{\varphi(d_-(k))}{\sqrt{w(k)}}.
\]

\end{proof}


Positivity of $q$, i.e.\ $g\ge0$, is however not the whole butterfly story: convexity of the
call slice controls its \emph{interior}, but the price must also decay at extreme strikes. The
complete characterisation is the following.

\newpage

\begin{lemma}[{Butterfly characterisation}]
\label{lem:gjbutterfly}
A $C^2$ slice $w(\cdot)>0$ is free of butterfly arbitrage if and only if
\begin{enumerate}[label=(\roman*)]
\item $g(k)\ge0$ for all $k\in\R$, and
\item $\displaystyle\lim_{k\to+\infty}d_+\bigl(k,w(k)\bigr)=-\infty$
(equivalently, $c(k,w(k))\to0$ as $k\to+\infty$).
\end{enumerate}
\end{lemma}

\begin{proof}[Proof]
By Proposition~\ref{prop:density}, $g\ge0$ is equivalent to the convexity of the call slice in
strike, i.e.\ absence of interior butterfly arbitrage.

For (ii): using \eqref{eq:blacknorm} and \eqref{eq:phiid} with the Mill's-ratio asymptotic
$N(x)\sim\varphi(x)/|x|$ as $x\to-\infty$,
\[
c=N(d_+)-e^kN(d_-)\quad\text{with}\quad
e^kN(d_-)\sim e^k\frac{\varphi(d_-)}{|d_-|}=\frac{\varphi(d_+)}{|d_-|}
\]
whenever $d_-\to-\infty$. If $d_+\to-\infty$ then both $N(d_+)\to0$ and
$\varphi(d_+)/|d_-|\to0$, so $c\to0$: the total mass of $q$ is $1$ and no mass escapes to
$+\infty$. Conversely, if $d_+\not\to-\infty$ along some subsequence then $N(d_+)$ stays
bounded away from $0$ faster than $e^kN(d_-)$ compensates and $c(k)$ does not vanish; a
nonvanishing call price at infinite strike is the classical arbitrage of selling calls of
arbitrarily large strike at a price bounded away from zero (see \cite{roper2010} and
\cite[\S2]{gatheraljacquier2014} for the complete argument, including the local-martingale
subtleties of \cite{coxhobson2005}).
\end{proof}

\begin{remark}[Lee's bound]\label{rem:lee}
Condition (ii) is a constraint on the \emph{wings}: by Lee's moment formula \cite{lee2004},
any arbitrage-free slice has $w(k)/|k|$ bounded by $2$ as $k\to\pm\infty$, so total variance
grows at most linearly with asymptotic slope at most $2$. The boundary case
$w(k)/|k|\to2$ is exactly $d_+\not\to-\infty$, i.e.\ the borderline of
Lemma~\ref{lem:gjbutterfly}(ii). The corresponding SVI wing condition
$b(1+|\rho|)\le2$ is derived in Section~\ref{sec:svisec}.
\end{remark}

% ============================================================================
\section{Calendar arbitrage and monotone total variance}\label{sec:calendarsec}
% ============================================================================

Butterfly-freeness is a per-slice property, and it can hold at every single maturity while the
family of slices remains mutually inconsistent. Even if each maturity is individually a valid
distribution, the slices must be correctly ordered relative to one another, since an option
cannot become cheaper by being granted more time. This section shows that the requirement is
exactly monotonicity of total variance in maturity.

One structural detail deserves attention here rather than later, because the whole of
Chapter~\ref{ch:learning} turns on it: the calendar condition involves a \emph{first} derivative
in $\tau$, whereas the butterfly condition involves a \emph{second} derivative in $k$. That
asymmetry looks like a technicality at this point. Section~\ref{sec:softtheory} shows it is the
mechanism behind the central finding of the thesis, because a constraint that can fail anywhere
along an axis and one that concentrates at one end of it require incompatible sampling grids.

\begin{definition}[Calendar arbitrage]\label{def:calendar}
The surface is free of calendar arbitrage if and only if
$\tau\mapsto w(k,\tau)$ is non-decreasing for every $k$. Equivalently, total-variance
slices do not cross.
\end{definition}


\begin{proposition}[Calendar characterisation]\label{prop:calendar}
Let $c(k,\tau):=c\bigl(k,w(k,\tau)\bigr)$ denote the normalised call surface. The following are
equivalent:
\begin{enumerate}[label=(\alph*)]
\item $\tau\mapsto w(k,\tau)$ is non-decreasing for every $k$;
\item $\tau\mapsto c(k,\tau)$ is non-decreasing for every $k$;
\item the surface is consistent, in the calendar direction, with \emph{some} nonnegative
martingale for the forward-normalised underlying, i.e.\ no calendar-spread portfolio (long the
longer-dated call, short the shorter-dated call, at equal moneyness $k$) has negative cost and
nonnegative value.
\end{enumerate}
\end{proposition}

\begin{proof}
\begin{enumerate}[label=\textup{(\roman*)}, leftmargin=*]

\item \textbf{Proof of $(a)\iff(b)$.}
By Lemma~\ref{lem:blackpartials}, $c_w>0$, so at fixed $k$ the Black price is strictly
increasing in total variance. Hence $\tau\mapsto w(k,\tau)$ is non-decreasing if and only if
$\tau\mapsto c(k,\tau)$ is non-decreasing.

\item \textbf{Proof of $(c)\Rightarrow(b)$.}
Let $(M_u)$ be a nonnegative martingale with $M_t=1$ representing the forward-normalised
underlying, so that
$c(k,\tau_i)=\E\bigl[(M_{T_i}-e^k)^+\bigr].$
For $\tau_1<\tau_2$, conditional Jensen's inequality applied to the convex function
$x\mapsto(x-e^k)^+$ gives
\[
\begin{aligned}
\E\bigl[(M_{T_2}-e^k)^+\bigr]
&=\E\Bigl[\E\bigl[(M_{T_2}-e^k)^+\mid\mathcal F_{T_1}\bigr]\Bigr]\\
&\ge \E\Bigl[\bigl(\E[M_{T_2}\mid\mathcal F_{T_1}]-e^k\bigr)^+\Bigr]\\
&=\E\bigl[(M_{T_1}-e^k)^+\bigr].
\end{aligned}
\]
Therefore, $c(k,\tau_2)\ge c(k,\tau_1)$.

\item \textbf{Proof of $(b)\Rightarrow(c)$.}
Assume that (b) holds and that every slice is free of butterfly arbitrage. The marginal laws
implied by the slices are then increasing in convex order. Kellerer's theorem
\cite{kellerer1972} yields a martingale with these marginals, so no static portfolio can
arbitrage the surface.

Conversely, suppose that (b) fails for some $\tau_1<\tau_2$ and $k$, so that $c(k,\tau_2)-c(k,\tau_1)<0.$
Consider a long $T_2$-call and a short $T_1$-call with the same forward moneyness $k$. The
portfolio has strictly negative initial cost. At $T_1$, the short leg pays
$-(M_{T_1}-e^k)^+$, while the value of the long leg is at least its intrinsic value by
Proposition~\ref{prop:shape}(b). The portfolio therefore has nonnegative value at $T_1$ and
defines a static arbitrage. The comparison uses equal moneyness rather than equal strike; in
discounted forward units, the strikes coincide.
\end{enumerate}
\end{proof}

\begin{remark}
Proposition~\ref{prop:calendar} is why the calendar penalty of every learning method in this
thesis is a penalty on $\partial_\tau w$: it is the complete, coordinate-free content of the
calendar constraint, and unlike the butterfly constraint, it involves only a \emph{first}
derivative. This asymmetry between the two constraints (first- vs second-order; anywhere on
the axis vs concentrated at the short end) is developed in Section~\ref{sec:softtheory} and is
the mechanism of the collocation blind spot.
\end{remark}

% ============================================================================
\section{Dupire's formula in $(k,w)$: the identity that prices the audit}\label{sec:dupiresec}
% ============================================================================

The two previous sections produced two audit quantities, $g$ and $\partial_\tau w$, and so far
they are only flags: a surface either satisfies them or does not. This section explains why they
are considerably more than flags, and it is the reason the audit of this thesis can be given an
economic denomination at all.

Dupire's insight is that an arbitrage-free surface of European prices already determines a
diffusion that reprices every vanilla exactly, and that the diffusion's local variance is a
ratio of derivatives of the surface. Written in our coordinates, that ratio turns out to be
$\partial_\tau w$ over $g$: the calendar quantity divided by the butterfly quantity. In other
words, the two things we audit are not incidental diagnostics but the numerator and the
denominator of the object a desk actually simulates with. A calendar violation is a negative
local variance; a butterfly violation is a vanishing or exploding one. We derive the equation
first, then change variables.

\begin{proposition}[Dupire's equation]\label{prop:dupirepde}
Work under zero rates and dividends (so $D\equiv1$, $F\equiv S_t$, and the forward measure
coincides with the risk-neutral one). Suppose the underlying follows the local-volatility
dynamics $\dd S_u=\sigma_{\mathrm{loc}}(u,S_u)\,S_u\,\dd W_u$ and that the transition density
$p(T,K)$ of $S_T$ is $C^{1,2}$ with $\sigma_{\mathrm{loc}}^2K^2p$ and its $K$-derivative
vanishing as $K\to\{0,\infty\}$. Then the undiscounted call $C(K,T)$ satisfies
\begin{equation}\label{eq:dupirepde}
\partial_T C(K,T)=\tfrac12\,\sigma_{\mathrm{loc}}^2(T,K)\,K^2\,\partial_{KK}C(K,T).
\end{equation}
\end{proposition}

\begin{proof}
The Kolmogorov forward (Fokker--Planck) equation for the density of the driftless diffusion is
\[
\partial_T p(T,K)=\tfrac12\,\partial_{KK}\bigl[\sigma_{\mathrm{loc}}^2(T,K)\,K^2\,p(T,K)\bigr].
\]
Write $C(K,T)=\int_0^\infty (s-K)^+p(T,s)\dd s=\int_K^\infty(s-K)p(T,s)\dd s$ and
differentiate under the integral:
\[
\partial_T C=\int_K^\infty(s-K)\,\partial_Tp(T,s)\dd s
=\tfrac12\int_K^\infty(s-K)\,\partial_{ss}\bigl[\sigma_{\mathrm{loc}}^2s^2p\bigr]\dd s.
\]
Integrating by parts twice, with the boundary terms at $s=\infty$ vanishing by assumption and
those at $s=K$ vanishing because $(s-K)$ and its derivative are evaluated where $(s-K)=0$
(first integration) and the surviving term is picked up at $s=K$ (second integration):
\[
\int_K^\infty(s-K)\,\partial_{ss}\bigl[\sigma^2s^2p\bigr]\dd s
=\Bigl[(s-K)\,\partial_s(\sigma^2s^2p)\Bigr]_K^\infty
-\int_K^\infty\partial_s(\sigma^2s^2p)\dd s
=0-\Bigl[\sigma^2s^2p\Bigr]_K^\infty
=\sigma_{\mathrm{loc}}^2(T,K)K^2p(T,K).
\]
Since $p(T,K)=\partial_{KK}C(K,T)$ (Proposition~\ref{prop:shape}(c)), \eqref{eq:dupirepde}
follows. Dupire's insight \cite{dupire1994} is to read \eqref{eq:dupirepde} backwards: given
an arbitrage-free call surface, \emph{define}
$\sigma^2_{\mathrm{loc}}=2\,\partial_TC/(K^2\partial_{KK}C)$; the resulting diffusion reprices
every vanilla exactly.
\end{proof}

\begin{proposition}[Local variance as calendar over butterfly]\label{prop:dupire}
Under the assumptions of Proposition~\ref{prop:dupirepde}, if the implied total-variance
surface $w(k,\tau)$ is $C^{2,1}$ with $g>0$, the Dupire local variance at $(k,\tau)$ is
\begin{equation}\label{eq:dupirekw}
\sigma_{\mathrm{loc}}^{2}(k,\tau)\;=\;\frac{\partial_\tau w(k,\tau)}{g(k,\tau)} .
\end{equation}
\end{proposition}

\begin{proof}
Substitute $C=F\,c\bigl(k,w(k,T)\bigr)$ into \eqref{eq:dupirepde}, with $k=\log(K/F)$ fixed as
$T$ varies (zero carry keeps $F$ constant).

\emph{Left side.} At fixed $K$ (hence fixed $k$), the only $T$-dependence is through $w$:
$\partial_TC=F\,c_w\,\partial_\tau w$, with $c_w=e^{k}\varphi(d_-)/(2\sqrt w)$ by
\eqref{eq:cw}. (The Black ``theta'' in these coordinates is carried entirely by $w$.)

\emph{Right side.} By \eqref{eq:qk} and Proposition~\ref{prop:density},
\[
K^2\,\partial_{KK}C
=F\bigl(\partial_{kk}c-\partial_kc\bigr)
=F\,e^{k}\,q(k)
=F\,e^{k}\,\frac{g\,\varphi(d_-)}{\sqrt w}.
\]
Equating,
\[
F\,\frac{e^{k}\varphi(d_-)}{2\sqrt w}\,\partial_\tau w
=\tfrac12\,\sigma_{\mathrm{loc}}^2\,F\,e^{k}\,\frac{g\,\varphi(d_-)}{\sqrt w}
\quad\Longleftrightarrow\quad
\partial_\tau w=\sigma_{\mathrm{loc}}^2\,g,
\]
which is \eqref{eq:dupirekw}. This is Gatheral's equation (1.10) \cite{gatheral2011book} in our
notation.
\end{proof}

\begin{remark}[Why this identity organises the thesis]\label{rem:organise}
Equation~\eqref{eq:dupirekw} says that the two no-arbitrage constraints are not abstract flags:
the \emph{calendar} constraint is the numerator of the local variance and the \emph{butterfly}
constraint is its denominator. A calendar violation ($\partial_\tau w<0$) is a negative local
variance; a butterfly violation ($g\le 0$) is a non-positive or exploding one. 
\end{remark}
